% Step-by-step derivation that accompanies theory-note.tex.
% Draft of 2026-11-19. Status and open items: see README.md in this folder.
\documentclass[aps,pra,onecolumn,notitlepage,11pt,amsmath,amssymb,nofootinbib]{revtex4-2}

\usepackage{bm}
\usepackage[colorlinks=true,allcolors=blue]{hyperref}

\newcommand{\Nop}{\hat N}
\newcommand{\Iop}{\hat I}
\newcommand{\dd}{\mathrm{d}}
\newcommand{\ee}{\mathrm{e}}
\newcommand{\ii}{\mathrm{i}}
\newcommand{\sinc}{\operatorname{sinc}}
\newcommand{\sgn}{\operatorname{sgn}}

\begin{document}

\title{Optical rectification with a quantized pump:\\
derivation of the exact solution from the Hamiltonian}
\author{Working notes to accompany the theory note, 19 November 2026
(cross-references updated 20 November)}
\noaffiliation

\begin{abstract}
These notes derive, step by step, the results of Sec.~II of the theory note:
the interaction Hamiltonian with an explicit coupling constant, the two
commutators that make it solvable, the input-output relation for the THz
field, the full evolution operator, the transformation of the pump, and the
THz state. Each step is marked as an identity or as an approximation. The
last section lists what was checked numerically and what was not.
\end{abstract}

\maketitle

\section{Result}

For a coupling that does not depend on the optical frequency across the pump
band, and without pump dispersion in the crystal, the evolution operator of
the pump and THz fields is
\begin{equation}
\hat U=\ee^{\ii\hat\Phi}\,\hat D,\qquad
\hat D=\exp\!\Big\{-\ii\!\int_0^\infty\!\!\dd\Omega\,
\big[\zeta(\Omega)\,\hat b^\dagger(\Omega)\Nop(\Omega)+\mathrm{H.c.}\big]\Big\},
\label{eq:result}
\end{equation}
with $\Nop(\Omega)$ the Fourier transform of the pump photon-flux operator and
$\hat\Phi$ a quadratic form in the flux. It follows that
\begin{equation}
\hat b_{\rm out}(\Omega)=\hat b_{\rm in}(\Omega)-\ii\zeta(\Omega)\Nop(\Omega),
\qquad \Iop_{\rm out}(\tau)=\Iop_{\rm in}(\tau),
\qquad \hat a_{\rm out}(\tau)=\ee^{-\ii\hat\vartheta(\tau)}\hat a_{\rm in}(\tau).
\label{eq:io}
\end{equation}
No expansion in the coupling strength is made.

\section{Fields}

\emph{Geometry.} One dimension along $z$, a crystal from $z=0$ to $z=L$, one
transverse mode of effective area $A$. Reflections at the faces are ignored.

\emph{THz field.} With continuum operators
$[\hat b(\Omega),\hat b^\dagger(\Omega')]=\delta(\Omega-\Omega')$, $\Omega>0$,
the forward-propagating field in the crystal is \cite{blow1990}
\begin{equation}
\hat E_T(z,t)=-\int_0^\infty\!\!\dd\Omega\;s(\Omega)
\Big[\hat b(\Omega)\,\ee^{\ii[k_T(\Omega)z-\Omega t]}+\mathrm{H.c.}\Big],
\qquad
s(\Omega)=\sqrt{\frac{\hbar\Omega}{4\pi\varepsilon_0c\,n_T(\Omega)A}} ,
\label{eq:ET}
\end{equation}
with $k_T=n_T\Omega/c$. The normalization is such that the energy carried
through a plane of area $A$ is
$\int\dd\Omega\,\hbar\Omega\,\hat b^\dagger(\Omega)\hat b(\Omega)$. The overall
factor $-1$ is a choice of phase for $\hat b$ (it
differs from the common convention by $\hat b\to-\ii\hat b$) and makes the
coupling in Eq.~\eqref{eq:Hint} real and positive for $\chi^{(2)}>0$.

\emph{Pump field.} The positive-frequency part of the optical field is
written with an envelope operator in the frame that moves with the group
velocity $v_g$,
\begin{equation}
\hat E_o^{(+)}(z,t)=\ii\sqrt{\frac{\hbar\omega_0}{2\varepsilon_0c\,n_oA}}\;
\hat a(\tau)\,\ee^{\ii(k_0z-\omega_0t)},\qquad \tau=t-\frac{z}{v_g},
\label{eq:Eo}
\end{equation}
\begin{equation}
\hat a(\tau)=\frac{1}{\sqrt{2\pi}}\int\!\dd\omega\,\hat a(\omega)\,
\ee^{-\ii(\omega-\omega_0)\tau},\qquad
[\hat a(\tau),\hat a^\dagger(\tau')]=\delta(\tau-\tau').
\end{equation}
Two approximations enter here. (A1) The prefactor is evaluated at the carrier
frequency $\omega_0$ (narrow relative bandwidth). (A2) The optical wave number
is expanded to first order, $k(\omega)=k_0+(\omega-\omega_0)/v_g$: no
group-velocity dispersion of the pump inside the crystal. With (A2) the
envelope depends on $z$ and $t$ only through $\tau$.

\emph{Photon flux.} The intensity is
$2\varepsilon_0cn_o\hat E_o^{(-)}\hat E_o^{(+)}=\hbar\omega_0\Iop(\tau)/A$ with
\begin{equation}
\Iop(\tau)=\hat a^\dagger(\tau)\hat a(\tau),\qquad
\Nop(\Omega)=\int\!\dd\tau\,\ee^{\ii\Omega\tau}\Iop(\tau)
=\int\!\dd\omega\,\hat a^\dagger(\omega)\hat a(\omega+\Omega),\qquad
\Nop^\dagger(\Omega)=\Nop(-\Omega).
\label{eq:N}
\end{equation}
$\Iop$ is in photons per second, $\Nop$ is dimensionless, and $\Nop(0)$ is the
photon number of the pulse.

\section{Interaction Hamiltonian}

\emph{Step 1: energy density.} To first order in $\chi^{(2)}$ the interaction
energy density of a medium with polarization
$P^{(2)}=\varepsilon_0\chi^{(2)}E^2$ is
\begin{equation}
u_{\rm NL}=-\tfrac13\,\varepsilon_0\chi^{(2)}E^3 .
\end{equation}
This is the form obtained when the displacement field is the canonical
variable \cite{quesada2017}; the coefficient and sign are confirmed in
Sec.~\ref{sec:classical} by comparison with the classical wave equation.

\emph{Step 2: select the rectification term.} With $E=E_o+E_T$,
\begin{equation}
E^3=E_o^3+3E_o^2E_T+3E_oE_T^2+E_T^3 .
\end{equation}
Only $3E_o^2E_T$ is kept (A3): $E_o^3$ does not involve the THz field, and the
last two terms couple THz modes among themselves and are of higher order in
the weak THz field. Hence $u_{\rm int}=-\varepsilon_0\chi^{(2)}E_o^2E_T$.

\emph{Step 3: rotating-wave approximation on the optical carrier.} With
$E_o=E_o^{(+)}+E_o^{(-)}$,
\begin{equation}
E_o^2=2\hat E_o^{(-)}\hat E_o^{(+)}+\big[\hat E_o^{(+)}\big]^2
+\big[\hat E_o^{(-)}\big]^2+\text{const}.
\end{equation}
The squares oscillate at $2\omega_0$ and cannot drive a THz field (A4). The
constant comes from normal ordering, does not depend on $\tau$, and has no
component at $\Omega\neq0$. What remains is the rectified polarization
\begin{equation}
\hat P_{\rm OR}(z,t)=2\varepsilon_0\chi^{(2)}\hat E_o^{(-)}\hat E_o^{(+)}
=\frac{\chi^{(2)}\hbar\omega_0}{n_oc\,A}\;\Iop\!\left(t-\frac{z}{v_g}\right),
\qquad u_{\rm int}=-\hat P_{\rm OR}\hat E_T ,
\label{eq:POR}
\end{equation}
a polarization proportional to the photon flux. The second form uses
Eq.~\eqref{eq:Eo}.

\emph{Step 4: integrate over the crystal.} Inserting Eq.~\eqref{eq:ET},
$\hat H_{\rm int}(t)=A\int_0^L\dd z\,u_{\rm int}$ becomes
\begin{equation}
\boxed{\;
\hat H_{\rm int}(t)=\hbar\int_0^L\!\!\dd z\int_0^\infty\!\!\dd\Omega\;
\kappa(\Omega)\,\Iop\!\left(t-\frac{z}{v_g}\right)
\Big[\hat b(\Omega)\,\ee^{\ii(k_Tz-\Omega t)}+\mathrm{H.c.}\Big]\;}
\label{eq:Hint}
\end{equation}
\begin{equation}
\boxed{\;\kappa(\Omega)=\frac{\chi^{(2)}\omega_0}{n_oc}
\sqrt{\frac{\hbar\Omega}{4\pi\varepsilon_0c\,n_T(\Omega)A}}\;}
\label{eq:kappa}
\end{equation}
This is the Hamiltonian in the interaction picture: the free propagation of
both fields is in the exponentials and in the argument of $\Iop$. The coupling
$\kappa$ has units $\mathrm{s}^{1/2}\,\mathrm{m}^{-1}$. Here
$\chi^{(2)}=\chi^{(2)}(\Omega;\omega_0,-\omega_0)$, related to the
electro-optic coefficient by $|\chi^{(2)}|=n_o^4r/2$; it is taken independent
of the optical frequency across the pump band (A5). Assumptions (A2) and (A5)
together are the ``flat coupling'' of the note.

\emph{Relation to the generator of the propagation picture.} Integrating
Eq.~\eqref{eq:Hint} over time with
$\int\dd t\,\Iop(t-z/v_g)\,\ee^{-\ii\Omega t}
=\ee^{-\ii\Omega z/v_g}\Nop^\dagger(\Omega)$ gives
\begin{equation}
\int_{-\infty}^{\infty}\!\!\dd t\,\hat H_{\rm int}(t)=\int_0^L\!\!\dd z\,\hat G(z),
\qquad
\hat G(z)=\hbar\int_0^\infty\!\!\dd\Omega\,\kappa(\Omega)
\Big[\ee^{-\ii\Delta k\,z}\,\hat b^\dagger(\Omega)\Nop(\Omega)+\mathrm{H.c.}\Big],
\label{eq:G}
\end{equation}
with $\Delta k(\Omega)=k_T(\Omega)-\Omega/v_g$. Equation~\eqref{eq:Hint} is
Eq.~(3) of the theory note, and $\hat G(z)$ is Eq.~(A6) of its Appendix~A
(numbers of the version of 20 November 2026). The time integral selects the energy-conserving combinations
$\hat b^\dagger\Nop$ (a pump photon is shifted down by $\Omega$ and a THz
photon is created) and $\hat b\,\Nop^\dagger$; no rotating-wave approximation
is needed for the THz field.

\section{The two commutators}

\emph{(C1) The flux commutes with itself.} Using
$\hat a(\tau)\hat a^\dagger(\tau')=\hat a^\dagger(\tau')\hat a(\tau)
+\delta(\tau-\tau')$,
\begin{equation}
\Iop(\tau)\Iop(\tau')=\hat a^\dagger(\tau)\hat a^\dagger(\tau')\hat a(\tau)\hat a(\tau')
+\delta(\tau-\tau')\,\hat a^\dagger(\tau)\hat a(\tau').
\end{equation}
The first term is symmetric under $\tau\leftrightarrow\tau'$, and so is the
second because of the delta function. Therefore
\begin{equation}
[\Iop(\tau),\Iop(\tau')]=0,\qquad
[\Nop(\Omega),\Nop(\Omega')]=[\Nop(\Omega),\Nop^\dagger(\Omega')]=0 .
\label{eq:C1}
\end{equation}
In the frequency domain the same statement reads
\begin{equation}
[\Nop(\Omega),\Nop^\dagger(\Omega')]=\int\!\dd\omega\,
\hat a^\dagger(\omega)\hat a(\omega+\Omega-\Omega')
-\int\!\dd\omega\,\hat a^\dagger(\omega+\Omega')\hat a(\omega+\Omega)=0,
\end{equation}
where the two integrals cancel after the shift
$\omega\to\omega-\Omega'$ in the second. The shift requires that the
integration range has no edge: this is where a finite band of flat coupling
would enter, and it is the origin of the edge term
$\hat n_1-\hat n_D$ in the numerical model of Sec.~\ref{sec:checks}.

\emph{(C2) The THz field has a c-number commutator.} With
$\hat\Theta(z,t)=\int\dd\Omega\,\kappa\,[\hat b\,\ee^{\ii(k_Tz-\Omega t)}
+\mathrm{H.c.}]$,
\begin{equation}
[\hat\Theta(z,t),\hat\Theta(z',t')]=2\ii\,K(z-z',t-t'),\qquad
K(\Delta z,\Delta t)=\int_0^\infty\!\!\dd\Omega\,\kappa^2(\Omega)
\sin\!\big[k_T\Delta z-\Omega\Delta t\big].
\label{eq:C2}
\end{equation}

\emph{Consequences.} The Hamiltonian is
$\hat H_{\rm int}(t)=\hbar\int\dd z\,\Iop(t-z/v_g)\hat\Theta(z,t)$, a product
of a pump operator and a THz operator that commute with each other. From (C1),
\begin{equation}
[\hat H_{\rm int}(t),\Iop(\tau)]=0\quad\text{for all }t,\tau:
\end{equation}
the photon flux at every instant is a constant of the motion. From (C1) and
(C2),
\begin{equation}
[\hat H_{\rm int}(t),\hat H_{\rm int}(t')]=2\ii\hbar^2\!\int_0^L\!\!\dd z\!\int_0^L\!\!\dd z'\;
\Iop\!\left(t-\frac{z}{v_g}\right)\Iop\!\left(t'-\frac{z'}{v_g}\right)K(z-z',t-t'),
\label{eq:HH}
\end{equation}
which contains flux operators only and therefore commutes with
$\hat H_{\rm int}$ at any third time.

\section{Input-output relation for the THz field}
\label{sec:heis}

This is the shortest route to the central result. Let
$\hat b(\Omega,t)=\hat U^\dagger(t)\hat b(\Omega)\hat U(t)$. Then
$\partial_t\hat b(\Omega,t)=(\ii/\hbar)\hat U^\dagger[\hat H_{\rm int}(t),\hat
b(\Omega)]\hat U$, and with
$[\hat b^\dagger(\Omega'),\hat b(\Omega)]=-\delta(\Omega-\Omega')$
\begin{equation}
\partial_t\hat b(\Omega,t)=-\ii\kappa(\Omega)\int_0^L\!\!\dd z\;
\Iop\!\left(t-\frac{z}{v_g}\right)\ee^{-\ii(k_Tz-\Omega t)} .
\label{eq:dbdt}
\end{equation}
The flux on the right-hand side is the Heisenberg operator, but because the
flux is conserved it equals the input operator. The equation is therefore
linear with a source that does not evolve, and it integrates directly. With
$\int\dd t\,\Iop(t-z/v_g)\,\ee^{\ii\Omega t}=\ee^{\ii\Omega z/v_g}\Nop(\Omega)$,
\begin{equation}
\hat b_{\rm out}(\Omega)=\hat b_{\rm in}(\Omega)
-\ii\kappa(\Omega)\int_0^L\!\!\dd z\,\ee^{-\ii\Delta k\,z}\;\Nop(\Omega)
=\hat b_{\rm in}(\Omega)-\ii\zeta(\Omega)\Nop(\Omega),
\end{equation}
\begin{equation}
\boxed{\;\zeta(\Omega)=\kappa(\Omega)\,L\;\ee^{-\ii\Delta kL/2}\,
\sinc\!\left(\frac{\Delta kL}{2}\right),\qquad \sinc x=\frac{\sin x}{x}.\;}
\label{eq:zeta}
\end{equation}
In a perturbative treatment one would replace the flux in
Eq.~\eqref{eq:dbdt} by its input value as an approximation, valid to first
order. Here that replacement is exact, because of (C1).

\section{Evolution operator}

\subsection{Magnus expansion}

The time-ordered exponential
$\hat U=\mathcal T\exp[-(\ii/\hbar)\int\dd t\,\hat H_{\rm int}(t)]$ can be
written as $\exp(\hat\Omega_1+\hat\Omega_2+\dots)$ \cite{blanes2009} with
\begin{equation}
\hat\Omega_1=-\frac{\ii}{\hbar}\int\!\dd t\,\hat H_{\rm int}(t),\qquad
\hat\Omega_2=-\frac{1}{2\hbar^2}\int\!\dd t\!\int^{t}\!\!\dd t'\,
[\hat H_{\rm int}(t),\hat H_{\rm int}(t')],
\end{equation}
and higher terms that are nested commutators of these with
$\hat H_{\rm int}$. By Eq.~\eqref{eq:HH} all of them vanish, so the series
ends at the second term, and because $\hat\Omega_2$ commutes with
$\hat\Omega_1$ the exponential factorizes.

\emph{First term.} From Eq.~\eqref{eq:G} and
$\int_0^L\dd z\,\ee^{-\ii\Delta kz}=\zeta/\kappa$,
\begin{equation}
\hat\Omega_1=-\ii\int_0^\infty\!\!\dd\Omega\,
\big[\zeta\,\hat b^\dagger\Nop+\zeta^*\,\hat b\,\Nop^\dagger\big],
\qquad \hat D=\ee^{\hat\Omega_1}.
\end{equation}

\emph{Second term.} Inserting Eq.~\eqref{eq:HH}, changing to
$\tau=t-z/v_g$, $\tau'=t'-z'/v_g$, and using that the integrand is
antisymmetric under the exchange of the two events,
\begin{equation}
\hat\Omega_2=\ii\hat\Phi,\qquad
\hat\Phi=-\iint\!\dd\tau\,\dd\tau'\;W(\tau-\tau')\,\Iop(\tau)\Iop(\tau'),
\label{eq:Phi}
\end{equation}
\begin{equation}
W(x)=\frac12\int_0^L\!\!\dd z\!\int_0^L\!\!\dd z'\;
\sgn\!\Big(x+\frac{z-z'}{v_g}\Big)
\int_0^\infty\!\!\dd\Omega\,\kappa^2\sin\!\big[\Delta k\,(z-z')-\Omega x\big].
\label{eq:W}
\end{equation}
Together: $\hat U=\ee^{\ii\hat\Phi}\hat D$, Eq.~\eqref{eq:result}.

\subsection{Ordering in $z$ and the kernel of $\hat\Phi$}

Forward propagation is often described by ordering the evolution in $z$ with
the generator $\hat G(z)$, $\partial_z\hat U=-(\ii/\hbar)\hat G(z)\hat U$; the
theory note gives this form in its Appendix~A. The first Magnus term is the same,
$\int\dd z\,\hat G=\int\dd t\,\hat H_{\rm int}$. In the second term the sign
function in Eq.~\eqref{eq:W} is replaced by $\sgn(z-z')$, and the integral can
then be done:
\begin{equation}
V(x)=\int_0^\infty\!\!\dd\Omega\;\varphi(\Omega)\cos\Omega x,\qquad
\varphi(\Omega)=\kappa^2\,\frac{\Delta kL-\sin\Delta kL}{\Delta k^2},
\label{eq:V}
\end{equation}
using $\int_0^L\dd z\int_0^z\dd z'\sin[\Delta k(z-z')]
=(\Delta kL-\sin\Delta kL)/\Delta k^2$. The two kernels differ only where the
later of two events lies upstream of the earlier one. A forward-propagating
THz wave gives no contribution there, except at coincidence. For a
dispersionless THz index and $\kappa^2=c_\kappa\Omega$ this can be made
explicit: $K=-\pi c_\kappa\,\delta'(n_T\Delta z/c-\Delta t)$ is supported on
the THz light cone, where the two sign functions agree, and
\begin{equation}
W(x)-V(x)=-\pi c_\kappa\,\frac{cL}{n_T}\,\delta(x).
\end{equation}
The difference is a term $\propto L\int\dd\tau\,\Iop(\tau)^2$ in $\hat\Phi$:
an instantaneous Kerr-type phase, of the kind known as the cascaded
contribution of rectification and the electro-optic effect to the nonlinear
index \cite{caumes2002}. Its coefficient depends on parts of the THz response
that a forward-wave model does not contain, so the model does not fix it. It
affects only the phase of the pump, Eq.~\eqref{eq:theta}. \emph{Everything
that concerns the THz field, and the conservation of the flux, is independent
of this choice.}

\subsection{The same result from Glauber's classical-source theorem}

Because the flux is conserved, the Hamiltonian is block diagonal in the joint
eigenstates of $\Iop(\tau)$. These are states with photons at definite times,
$|\tau_1,\dots,\tau_n\rangle\propto\prod_i\hat a^\dagger(\tau_i)|0\rangle$,
with eigenvalue $I(\tau)=\sum_i\delta(\tau-\tau_i)$. In each block
\begin{equation}
\hat H_{\rm int}(t)=\hbar\int\!\dd\Omega\,\big[j^*(\Omega,t)\,\hat b(\Omega)
+j(\Omega,t)\,\hat b^\dagger(\Omega)\big],\qquad
j(\Omega,t)=\kappa\int_0^L\!\!\dd z\,I\!\left(t-\frac{z}{v_g}\right)
\ee^{-\ii(k_Tz-\Omega t)},
\end{equation}
a set of harmonic oscillators driven by a prescribed c-number source. The
evolution operator of a driven oscillator is a displacement times a phase
\cite{glauber1963}, with amplitude
\begin{equation}
\beta(\Omega)=-\ii\int\!\dd t\,j(\Omega,t)
=-\ii\,\zeta(\Omega)\sum_i\ee^{\ii\Omega\tau_i}.
\label{eq:beta}
\end{equation}
This is Eq.~\eqref{eq:result} evaluated in one block: $\Nop(\Omega)$ has the
eigenvalue $\sum_i\ee^{\ii\Omega\tau_i}$ and $\hat\Phi$ is a number.

\section{Transformation of the pump}

Write $\hat D=\exp[-\ii\int\dd\tau'\,\Iop(\tau')\hat F_{\rm in}(\tau')]$ with
$\hat F(\tau)=\int\dd\Omega\,[\zeta\,\hat b^\dagger\ee^{\ii\Omega\tau}
+\mathrm{H.c.}]$, and use
\begin{equation}
[\Iop(\tau'),\hat a(\tau)]=-\delta(\tau-\tau')\,\hat a(\tau),\qquad
[\hat F(\tau'),\hat F(\tau)]=-2\ii S(\tau'-\tau),\quad
S(x)=\int_0^\infty\!\!\dd\Omega\,|\zeta|^2\sin\Omega x .
\end{equation}

\emph{Step 1: conjugation with $\hat D$.} Let
$\hat A(\lambda)=\hat D^{-\lambda}\hat a(\tau)\hat D^{\lambda}$. Then
$\partial_\lambda\hat A=-\ii\hat F_\lambda(\tau)\hat A$ with
$\hat F_\lambda(\tau)=\hat D^{-\lambda}\hat F_{\rm in}(\tau)\hat D^{\lambda}
=\hat F_{\rm in}(\tau)+2\lambda\int\dd\tau'\,S(\tau'-\tau)\Iop(\tau')$. The
operators $\hat F_\lambda$ at different $\lambda$ commute, and $S(0)=0$ removes
any ordering question with $\hat a(\tau)$, so
\begin{equation}
\hat D^\dagger\hat a(\tau)\hat D=\exp\!\Big\{-\ii\hat F_{\rm in}(\tau)
-\ii\!\int\!\dd\tau'\,S(\tau'-\tau)\Iop(\tau')\Big\}\hat a(\tau)
=\exp\!\Big\{-\tfrac{\ii}{2}\big[\hat F_{\rm in}(\tau)+\hat F_{\rm out}(\tau)\big]\Big\}\hat a(\tau).
\end{equation}
The second form uses $\hat F_{\rm out}-\hat F_{\rm in}
=2\int\dd\tau'\,S(\tau'-\tau)\Iop(\tau')$, which follows from
Eq.~\eqref{eq:io}.

\emph{Step 2: conjugation with $\ee^{\ii\hat\Phi}$.} From
Eq.~\eqref{eq:Phi},
$[\hat\Phi,\hat a(\tau)]=\hat\theta(\tau)\hat a(\tau)$ with
$\hat\theta(\tau)=2\int\dd\tau'\,W(\tau-\tau')\Iop(\tau')+W(0)$, and
$\ee^{-\ii\hat\Phi}\hat a(\tau)\ee^{\ii\hat\Phi}=\ee^{-\ii\hat\theta(\tau)}\hat a(\tau)$.

\emph{Result.}
\begin{equation}
\boxed{\;\hat a_{\rm out}(\tau)=\ee^{-\ii\hat\vartheta(\tau)}\hat a_{\rm in}(\tau),\qquad
\hat\vartheta(\tau)=\tfrac12\big[\hat F_{\rm in}(\tau)+\hat F_{\rm out}(\tau)\big]
+\hat\theta(\tau).\;}
\label{eq:theta}
\end{equation}
The pump is phase modulated and nothing else: this is the electro-optic
effect, with the THz field the pump travels with. Since $\hat\vartheta$ is
Hermitian, $\Iop_{\rm out}=\hat a_{\rm out}^\dagger\hat a_{\rm out}=\Iop_{\rm in}$,
which confirms the conservation of the flux.

\section{State of the THz field}

Expand the pump state in the photon-arrival basis and take vacuum for the THz
input. By Eq.~\eqref{eq:beta} each component becomes
$\ee^{\ii\Phi_{\{\tau\}}}|\tau_1,\dots,\tau_n\rangle\otimes|\beta_{\{\tau\}}\rangle$
with a multimode coherent state of the THz field. Tracing over the pump
removes the cross terms between different configurations, which are
orthogonal, and the phases:
\begin{equation}
\hat\rho_T=\sum_{n=0}^\infty\int\!\dd^n\tau\;p_n(\tau_1,\dots,\tau_n)\,
|\beta_{\{\tau\}}\rangle\langle\beta_{\{\tau\}}| ,
\end{equation}
where $p_n$ is the probability density of finding $n$ pump photons at times
$\tau_1<\dots<\tau_n$. In the time domain, with Eq.~\eqref{eq:ET} evaluated
at the exit face,
\begin{equation}
\hat E_T(L,t)=\hat E_{\rm vac}(t)+\int\!\dd\tau\,g(t-\tau)\,\Iop(\tau),\qquad
g(t)=2\,\mathrm{Re}\int_0^\infty\!\!\dd\Omega\;\ii\,s(\Omega)\zeta(\Omega)\,
\ee^{\ii(k_TL-\Omega t)} .
\label{eq:g}
\end{equation}
$g(t)$ is the field radiated per pump photon.

\section{Check against the classical result}
\label{sec:classical}

Classically, the forward THz amplitude
$\tilde E(z,\Omega)=E_T(z,\Omega)\ee^{-\ii k_Tz}$ obeys the slowly varying
wave equation
\begin{equation}
\partial_z\tilde E(z,\Omega)=\ii\,\frac{\Omega}{2\varepsilon_0n_Tc}\,
P_{\rm OR}(z,\Omega)\,\ee^{-\ii k_Tz},\qquad
P_{\rm OR}(z,t)=\frac{\chi^{(2)}}{n_oc}\,I_W\!\left(t-\frac{z}{v_g}\right),
\end{equation}
for an intensity $I_W$ in $\mathrm{W\,m^{-2}}$. Integration over the crystal
gives
\begin{equation}
\tilde E(L,\Omega)=\ii\,\frac{\Omega\,\chi^{(2)}L}{2\varepsilon_0n_Tn_oc^2}\;
\ee^{-\ii\Delta kL/2}\sinc\!\left(\frac{\Delta kL}{2}\right)\tilde I_W(\Omega),
\qquad \tilde I_W(\Omega)=\int\!\dd\tau\,\ee^{\ii\Omega\tau}I_W(\tau).
\end{equation}
The quantum mean field from Eqs.~\eqref{eq:ET}, \eqref{eq:io},
\eqref{eq:zeta}, and \eqref{eq:kappa} is
\begin{equation}
\langle\hat E_T^{(+)}(L,t)\rangle=\int_0^\infty\!\!\dd\Omega\;
\ii\,s\,\zeta\,\langle\Nop(\Omega)\rangle\,\ee^{\ii(k_TL-\Omega t)},\qquad
s\,\kappa L=\frac{1}{2\pi}\,\frac{\Omega\,\chi^{(2)}L}{2\varepsilon_0n_Tn_oc^2}\,
\frac{\hbar\omega_0}{A}.
\end{equation}
With $\tilde I_W=(\hbar\omega_0/A)\langle\Nop\rangle$ and
$E^{(+)}(t)=\int_0^\infty(\dd\Omega/2\pi)\,E(\Omega)\ee^{-\ii\Omega t}$ the two
expressions are identical, in magnitude and in phase. This fixes the
coefficient and the sign in Step 1 of Sec.~III, and it shows that the mean
THz field is the classical one for the mean intensity, whatever the state of
the pump.

Two quantities of use follow from $\kappa$. The number of THz photons is
$\int\dd\Omega\,|\zeta|^2\langle\Nop^\dagger\Nop\rangle$. The measurement
strength of the note is
\begin{equation}
B_T=\int_0^\infty\!\!\dd\Omega\,|\zeta(\Omega)|^2
=\left(\frac{\chi^{(2)}\omega_0L}{n_oc}\right)^{\!2}\frac{\hbar}{4\pi\varepsilon_0cA}
\int_0^\infty\!\!\dd\Omega\,\frac{\Omega}{n_T(\Omega)}\,
\sinc^2\!\left(\frac{\Delta kL}{2}\right),
\end{equation}
to be cut off by the bandwidth over which the THz field is collected.

\section{Assumptions and checks}
\label{sec:checks}

\emph{Assumptions}, in the order they entered: one transverse mode and no
reflections; (A1) narrow relative bandwidth in the field normalization; (A2)
no pump group-velocity dispersion; (A3) only the term linear in the THz field;
(A4) rotating-wave approximation on the optical carrier; (A5) $\chi^{(2)}$
independent of the optical frequency; no absorption of the pump. Linear
absorption and dispersion of the THz wave are compatible with the derivation:
they make $k_T$ complex and change $\zeta$ and $g$, and for an absorbing
medium the noise operators of the loss add vacuum fluctuations. Approximation
(A1) is needed for (C1): with the frequency-dependent prefactor kept, the
rectified polarization is no longer exactly proportional to the photon flux,
and (C1) acquires corrections of relative order $\Omega/\omega_0$, about one
percent for THz frequencies and a near-infrared pump.

\emph{Checked numerically.} The tests are in the repository under
\texttt{code/tests}, in the file for optical rectification. They use a model
with a ladder of $D$ optical frequency bins,
$\Nop=\sum_j\hat a_j^\dagger\hat a_{j+1}$, one THz mode, two pump photons, and
a full matrix exponential.
\begin{itemize}
\item Eq.~\eqref{eq:io} for the THz field and the conservation of the flux, at
couplings $|\zeta|=0.3$ and $0.6$: agreement to better than $10^{-6}$, limited
by the population of the end bins, where $[\Nop,\Nop^\dagger]=\hat n_1-\hat n_D$.
\item Eq.~\eqref{eq:theta} at perfect phase matching, through the pump
coherence $\langle\hat a_j^\dagger\hat a_k\rangle$: agreement to
$7\times10^{-10}$ and $2\times10^{-7}$; without the flux-dependent part of the phase the difference is
$0.03$ and $0.09$.
\item Eq.~\eqref{eq:result} with a phase mismatch ($\Delta kL=2.4$,
$|\zeta|=0.47$, $\varphi=0.11$), with $z$ ordering and the kernel
\eqref{eq:V}: the state $\ee^{\ii\hat\Phi}\hat D|\psi\rangle$ agrees with the
exact $z$-ordered evolution to $1\times10^{-4}$, $2\times10^{-5}$, and
$4\times10^{-6}$ in norm for $D=13$, $15$, and $17$; without $\hat\Phi$ the
difference is $0.30$.
\item The coefficient $\kappa$ against the classical formula: analytically,
Sec.~\ref{sec:classical}.
\end{itemize}

\emph{Not checked.} The derivation has not been repeated independently by a
second person. The time-ordered kernel \eqref{eq:W} and its relation to
\eqref{eq:V} are analytic only, and the statement about their difference is
shown for a dispersionless THz index only. The numerical model has one THz
mode, so it does not test the integrals over $\Omega$. Pump dispersion and
pump absorption, which break (C1), are not treated.

\section{Relation to known exact solutions}

The mechanism used here, a conserved photon flux coupled linearly to
harmonic oscillators, is not new. It is the basis of the exact quantum theory
of self-phase modulation without dispersion \cite{blow1991} and of its
extension to a response of finite duration, in which the field is coupled to
oscillators of the medium and the Heisenberg equation is solved in closed
form \cite{boivin1994}; of continuous-time models of cross-phase modulation,
where the phase of one beam is proportional to the photon flux of the other
and a phase-noise operator appears \cite{shapirorazavi2007}; and, for a
single mode, of the optomechanical and cavity electro-optic Hamiltonian
$\hbar g\hat a^\dagger\hat a(\hat b+\hat b^\dagger)$. Equation
\eqref{eq:theta} has the form reported for those models, a phase proportional
to the flux plus a phase-noise operator. What is specific to
optical rectification is that the oscillators are the modes of the emitted
THz field, so that the oscillator side of the solution,
Eqs.~\eqref{eq:io} and \eqref{eq:beta}, is the measured output. Only the
abstracts of Refs.~\cite{blow1991,boivin1994} were available to us; a
comparison with their full text remains to be made.

\begin{thebibliography}{9}
\bibitem{blow1990} K. J. Blow, R. Loudon, S. J. D. Phoenix, and T. J.
Shepherd, Continuum fields in quantum optics, Phys. Rev. A \textbf{42}, 4102
(1990) [entered from memory, to be checked].
\bibitem{quesada2017} N. Quesada and J. E. Sipe, Why you should not use the
electric field to quantize in nonlinear optics, Opt. Lett. \textbf{42}, 3443
(2017) [entered from memory, to be checked].
\bibitem{blanes2009} S. Blanes, F. Casas, J. A. Oteo, and J. Ros, The Magnus
expansion and some of its applications, Phys. Rep. \textbf{470}, 151 (2009)
[entered from memory, to be checked].
\bibitem{glauber1963} R. J. Glauber, Coherent and incoherent states of the
radiation field, Phys. Rev. \textbf{131}, 2766 (1963) [entered from memory,
to be checked].
\bibitem{caumes2002} J.-P. Caumes, L. Videau, C. Rouyer, and E. Freysz,
Kerr-like nonlinearity induced via terahertz generation and the
electro-optical effect in zinc blende crystals, Phys. Rev. Lett.
\textbf{89}, 047401 (2002) [entered from memory, to be checked].
\bibitem{blow1991} K. J. Blow, R. Loudon, and S. J. D. Phoenix, Exact
solution for quantum self-phase modulation, J. Opt. Soc. Am. B \textbf{8},
1750 (1991) [abstract read; full text not available].
\bibitem{boivin1994} L. Boivin, F. X. K\"artner, and H. A. Haus, Analytical
solution to the quantum field theory of self-phase modulation with a finite
response time, Phys. Rev. Lett. \textbf{73}, 240 (1994) [abstract read; full
text not available; author list from memory].
\bibitem{shapirorazavi2007} J. H. Shapiro and M. Razavi, Continuous-time
cross-phase modulation and quantum computation, New J. Phys. \textbf{9}, 16
(2007) [abstract read].
\end{thebibliography}

\end{document}
